% =============================================================================
% Custom preamble for simplicial_vector_calculus (pandoc -> XeLaTeX)
% Pulled in via `--include-in-header` by build.ps1 / build.sh.
%
% This file is injected into pandoc's default LaTeX template at the
% `header-includes` slot, which runs BEFORE pandoc's own \usepackage{hyperref}
% and \usepackage{unicode-math}. We therefore:
%   * load `mathtools` here (so it precedes unicode-math, per its manual),
%   * defer any \hypersetup{...} with \AtBeginDocument{...}.
% We only ADD or tune; we do not reload packages pandoc already loads
% (amsmath, amssymb, amsfonts, hyperref, xcolor, fontspec, unicode-math).
% =============================================================================

% ---- Math / symbols (mathtools must precede unicode-math) ------------------
\usepackage{mathtools}
\usepackage{amsthm}

% ---- Typography -------------------------------------------------------------
\usepackage{microtype}
\usepackage{booktabs}
% Let LaTeX stretch lines a bit further when it can't find a clean break;
% keeps overfull hboxes out of display math and URL-heavy lines.
\setlength{\emergencystretch}{3em}
% Make long URLs breakable at hyphens (applied before hyperref loads).
\PassOptionsToPackage{hyphens}{url}

% ---- Section heading polish -------------------------------------------------
\usepackage{titlesec}
\titleformat{\section}
  {\normalfont\Large\bfseries}{\thesection}{0.75em}{}
\titleformat{\subsection}
  {\normalfont\large\bfseries}{\thesubsection}{0.6em}{}
\titleformat{\subsubsection}
  {\normalfont\normalsize\bfseries}{\thesubsubsection}{0.5em}{}

% ---- Hyperref polish (hyperref is loaded later by pandoc; defer) ------------
% Colors are defined here so they're available wherever pandoc finally calls
% \hypersetup via its own template variables; the deferred \hypersetup below
% just adds PDF metadata on top of what pandoc's -V flags already set.
\AtBeginDocument{%
  \hypersetup{%
    pdftitle={Intrinsic Vector Calculus on Simplicial (Quadray) Coordinates},%
    pdfauthor={Leonardo Murillo Montero},%
    pdfsubject={Simplicial vector calculus; cross product; rotation; zero-sum hyperplane},%
    pdfkeywords={simplicial coordinates, Quadray coordinates, barycentric coordinates, gauge direction, zero-sum hyperplane, cross product, Hodge dual, wedge product, Rodrigues formula, Lie algebra, rotation}%
  }%
}

% ---- Theorem-like environments (available if source ever migrates) ---------
% The markdown source uses `**Theorem 4.1 (title).**` bold-header paragraphs
% rather than native environments, so these are currently unused by the body
% but are kept here so a future migration does not require a preamble change.
\theoremstyle{plain}
\newtheorem{theorem}{Theorem}[section]
\newtheorem{proposition}[theorem]{Proposition}
\newtheorem{lemma}[theorem]{Lemma}
\newtheorem{corollary}[theorem]{Corollary}
\theoremstyle{definition}
\newtheorem{definition}[theorem]{Definition}
\theoremstyle{remark}
\newtheorem{remark}[theorem]{Remark}
\newtheorem{example}[theorem]{Example}
